\section{Limitações}
\label{sec:limitacoes}

O DRM Transformer é um protótipo de pesquisa. O baseline atual é pequeno,
aproximadamente 3,5 milhões de parâmetros, e foi treinado por poucos tokens em
comparação com modelos de linguagem modernos. Portanto, os resultados de
perplexidade não devem ser comparados diretamente a modelos de produção.

A validação topológica depende de escolhas de avaliação: número de sementes
Voronoi, tamanho do subsample de homologia, razão de barra longa, projeção PCA,
filtro de densidade e número de restarts. O critério estrito
$f_{T^2}\geq 0{,}60$ reduz falsos positivos, mas ainda é uma convenção
experimental. Nos resultados atuais, o modelo treinado apresenta assinatura
toroidal quase-estável, mas não atinge esse limiar estrito.

As perdas toroidais são indutivas: elas ajudam a induzir $T^2$ em dimensão
quatro, mas também podem competir com a perda de linguagem. O fato de controles
sem regularização toroidal ainda apresentarem algum sinal próximo de $T^2$
sugere que a arquitetura também contribui para a estrutura, mas a separação
entre topologia aprendida, topologia induzida e artefatos de TDA ainda requer
mais ablações.

Além disso, a geometria local permanece espessa em 4D em várias rodadas. A
assinatura toroidal aparece com mais clareza em homologia PCA3 com filtro de
densidade, e não de forma plenamente estável no embedding 4D bruto. Em modelos
maiores, os pesos das perdas geométricas provavelmente precisarão ser
reescalados.

Por fim, a relação com a física relativística é formal e analógica. O fator de
Lorentz é usado como inspiração matemática para controle de escala; o modelo
não afirma implementar uma nova teoria física.

De modo semelhante, a discussão sobre geometria aberta, negociação e segurança
de IA é uma motivação conceitual, não uma garantia operacional. O DRM
Transformer não prova que uma IA será segura por possuir topologia fechada, nem
substitui avaliação comportamental, governança, interpretabilidade ou controle
externo. A contribuição aqui é propor uma hipótese mensurável: se certos
conflitos epistêmicos e semânticos aparecem como estruturas geométricas, então
podem ser comparados por ablações, homologia persistente e controles de
inicialização.
